% GENERATED FILE. Edit canonical lessons, then run npm run generate:books.
% canonical-source-hash: fnv1a64:f17a19477a6d1d23
% canonical-lessons: RU-C464-prognoz, RU-C464-sinoptik, RU-C464-osadki, RU-C464-pasmurnyy, RU-C464-proyasnyatsya

\chapter{The weather forecast}
\label{ch:ru-the-weather-forecast}
\hlchaptermodality{464}

\begin{chapteropening}
\textbf{By the end of this chapter:} \emph{I can follow a weather forecast: the forecast, the forecaster, rain and snow, an overcast sky and clearing up.}

Complete the last lesson of chapter 464: I can follow a weather forecast: the forecast, the forecaster, rain and snow, an overcast sky and clearing up.
\end{chapteropening}

\section[\texorpdfstring{\ru{прогноз} (prognóz)}{прогноз (prognóz)}]{\ru{прогноз} (prognóz) — a forecast}

\label{lesson:RU-C464-prognoz}



\begin{quote}
Before the new one: say the Russian for a seashell, then the Russian for to dive.
\end{quote}

\subsection*{You'll want to know: \ru{прогноз}}
\textbf{\ru{прогноз}} (\emph{prognóz}) — \textquotedblleft{}a forecast\textquotedblright{}.

\begin{quote}
\ru{По} \ru{прогнозу} \ru{завтра} \ru{будет} \ru{дождь}.
\end{quote}

\begin{quote}
\emph{Po prognózu závtra búdet dozhd.}
\end{quote}

\begin{quote}
According to the forecast it will rain tomorrow.
\end{quote}

Masculine. \emph{\ru{Прогноз} \ru{погоды}} — the weather forecast.

\begin{cousinweb}
\emph{\ru{прогноз}} is usually traced through Latin \emph{prognōsis} to Greek \emph{prógnōsis} 'knowing beforehand' (EtymDB lists both; the Russian entry has no source), from \emph{pro-} 'before' and \emph{gnōsis} 'knowledge'.
\end{cousinweb}

\begin{grammarlens}[title={What you've built}]
One more word for telling, explaining and describing at B1.
\end{grammarlens}

\subsection*{Your turn}
\begin{itemize}
\raggedright
  \item \emph{Say it:} \emph{prognóz}
  \item \emph{Say it:} \emph{prognóz}, once more
  \item \emph{Say it:} a sentence of your own with \emph{prognóz}
  \item \emph{Recall:} say the Russian for a seashell, then the Russian for to dive, then say \emph{prognóz} again
\end{itemize}

\begin{tcolorbox}[breakable,colback=teal!4,colframe=teal!35!black,title={Before you move on}]
What does \ru{прогноз} mean? (\textquotedblleft{}A forecast\textquotedblright{}.) Say it once more.
\end{tcolorbox}
\section[\texorpdfstring{\ru{синоптик} (sinóptik)}{синоптик (sinóptik)}]{\ru{синоптик} (sinóptik) — a weather forecaster}

\label{lesson:RU-C464-sinoptik}



\begin{quote}
Before the new one: say the Russian for to dive, then the Russian for a forecast.
\end{quote}

\subsection*{You'll want to know: \ru{синоптик}}
\textbf{\ru{синоптик}} (\emph{sinóptik}) — \textquotedblleft{}a weather forecaster\textquotedblright{}.

\begin{quote}
\ru{Синоптики} \ru{обещают} \ru{тёплую} \ru{неделю}.
\end{quote}

\begin{quote}
\emph{Sinóptiki obeshcháyut tyópluyu nedélyu.}
\end{quote}

\begin{quote}
The forecasters promise a warm week.
\end{quote}

Masculine, for a woman too.

\begin{cousinweb}
\emph{\ru{синоптик}} is usually traced to Greek \emph{synoptikós} 'seeing the whole together at a glance' (EtymDB lists the Greek word but not the Russian): a forecaster reads the whole weather map at once.
\end{cousinweb}

\begin{grammarlens}[title={What you've built}]
One more word for telling, explaining and describing at B1.
\end{grammarlens}

\subsection*{Your turn}
\begin{itemize}
\raggedright
  \item \emph{Say it:} \emph{sinóptik}
  \item \emph{Say it:} \emph{sinóptik}, once more
  \item \emph{Say it:} a sentence of your own with \emph{sinóptik}
  \item \emph{Recall:} say the Russian for to dive, then the Russian for a forecast, then say \emph{sinóptik} again
\end{itemize}

\begin{tcolorbox}[breakable,colback=teal!4,colframe=teal!35!black,title={Before you move on}]
What does \ru{синоптик} mean? (\textquotedblleft{}A weather forecaster\textquotedblright{}.) Say it once more.
\end{tcolorbox}
\section[\texorpdfstring{\ru{осадки} (osádki)}{осадки (osádki)}]{\ru{осадки} (osádki) — precipitation (rain or snow)}

\label{lesson:RU-C464-osadki}



\begin{quote}
Before the new one: say the Russian for a forecast, then the Russian for a weather forecaster.
\end{quote}

\subsection*{You'll want to know: \ru{осадки}}
\textbf{\ru{осадки}} (\emph{osádki}) — \textquotedblleft{}precipitation (rain or snow)\textquotedblright{}.

\begin{quote}
\ru{Сегодня} \ru{без} \ru{осадков}.
\end{quote}

\begin{quote}
\emph{Sevódnya bez osádkov.}
\end{quote}

\begin{quote}
Today there will be no rain or snow.
\end{quote}

Plural; a weather word. In the singular, \emph{\ru{осадок}} is a sediment.

\begin{cousinweb}
\emph{\ru{осадки}} is the plural of \emph{\ru{осадок}} 'a sediment, what settles', from \emph{\ru{о}-} and \emph{\ru{садить}} 'to set down'. EtymDB lists \emph{\ru{осадки}} without a deeper source.
\end{cousinweb}

\begin{grammarlens}[title={What you've built}]
One more word for telling, explaining and describing at B1.
\end{grammarlens}

\subsection*{Your turn}
\begin{itemize}
\raggedright
  \item \emph{Say it:} \emph{osádki}
  \item \emph{Say it:} \emph{osádki}, once more
  \item \emph{Say it:} a sentence of your own with \emph{osádki}
  \item \emph{Recall:} say the Russian for a forecast, then the Russian for a weather forecaster, then say \emph{osádki} again
\end{itemize}

\begin{tcolorbox}[breakable,colback=teal!4,colframe=teal!35!black,title={Before you move on}]
What does \ru{осадки} mean? (\textquotedblleft{}Precipitation (rain or snow)\textquotedblright{}.) Say it once more.
\end{tcolorbox}
\section[\texorpdfstring{\ru{пасмурный} (pásmurnyy)}{пасмурный (pásmurnyy)}]{\ru{пасмурный} (pásmurnyy) — overcast, cloudy, gloomy}

\label{lesson:RU-C464-pasmurnyy}



\begin{quote}
Before the new one: say the Russian for a weather forecaster, then the Russian for precipitation.
\end{quote}

\subsection*{You'll want to know: \ru{пасмурный}}
\textbf{\ru{пасмурный}} (\emph{pásmurnyy}) — \textquotedblleft{}overcast, cloudy, gloomy\textquotedblright{}.

\begin{quote}
\ru{Был} \ru{пасмурный} \ru{и} \ru{холодный} \ru{день}.
\end{quote}

\begin{quote}
\emph{Byl pásmurnyy i kholódnyy den.}
\end{quote}

\begin{quote}
It was an overcast, cold day.
\end{quote}

\emph{\ru{Пасмурно}} — it is overcast.

\begin{cousinweb}
\emph{\ru{пасмурный}} is usually linked with the old word \emph{\ru{смурый}} 'dark grey, gloomy'. EtymDB does not list it, so the line is hedged.
\end{cousinweb}

\begin{grammarlens}[title={What you've built}]
One more word for telling, explaining and describing at B1.
\end{grammarlens}

\subsection*{Your turn}
\begin{itemize}
\raggedright
  \item \emph{Say it:} \emph{pásmurnyy}
  \item \emph{Say it:} \emph{pásmurnyy}, once more
  \item \emph{Say it:} a sentence of your own with \emph{pásmurnyy}
  \item \emph{Recall:} say the Russian for a weather forecaster, then the Russian for precipitation, then say \emph{pásmurnyy} again
\end{itemize}

\begin{tcolorbox}[breakable,colback=teal!4,colframe=teal!35!black,title={Before you move on}]
What does \ru{пасмурный} mean? (\textquotedblleft{}Overcast, cloudy, gloomy\textquotedblright{}.) Say it once more.
\end{tcolorbox}
\section[\texorpdfstring{\ru{проясняться} (proyasnyátsya)}{проясняться (proyasnyátsya)}]{\ru{проясняться} (proyasnyátsya) — to clear up (of the weather)}

\label{lesson:RU-C464-proyasnyatsya}



\begin{quote}
Before the new one: say the Russian for precipitation, then the Russian for overcast.
\end{quote}

\subsection*{You'll want to know: \ru{проясняться}}
\textbf{\ru{проясняться}} (\emph{proyasnyátsya}) — \textquotedblleft{}to clear up (of the weather)\textquotedblright{}.

\begin{quote}
\ru{К} \ru{вечеру} \ru{небо} \ru{проясняется}.
\end{quote}

\begin{quote}
\emph{K vécheru nébo proyasnyáyetsya.}
\end{quote}

\begin{quote}
Towards evening the sky clears up.
\end{quote}

Imperfective; perfective \emph{\ru{проясниться}}. Built on \emph{\ru{ясный}} — clear.

\begin{cousinweb}
\emph{\ru{проясняться}} is built from \emph{\ru{про}-}, \emph{\ru{ясный}} 'clear, bright' and \emph{-\ru{ся}}. EtymDB does not list the verb.
\end{cousinweb}

\begin{grammarlens}[title={What you've built}]
One more word for telling, explaining and describing at B1.
\end{grammarlens}

\subsection*{Your turn}
\begin{itemize}
\raggedright
  \item \emph{Say it:} \emph{proyasnyátsya}
  \item \emph{Say it:} \emph{proyasnyátsya}, once more
  \item \emph{Say it:} a sentence of your own with \emph{proyasnyátsya}
  \item \emph{Recall:} say the Russian for precipitation, then the Russian for overcast, then say \emph{proyasnyátsya} again
\end{itemize}

\begin{tcolorbox}[breakable,colback=teal!4,colframe=teal!35!black,title={Before you move on}]
What does \ru{проясняться} mean? (\textquotedblleft{}To clear up (of the weather)\textquotedblright{}.) Say it once more.
\end{tcolorbox}
